\section{数值积分与独立复核}
给定控制后，式\eqref{eq:dynamics}构成非线性初值问题。搜索阶段使用双精度RK4、最大步长0.05\,s；取整步数后均匀调整步长，使末点与计划关机时刻一致。启停使用0.5\,s线性过渡。三算法共用积分器和同一预生成扰动样本。

最终控制保持不变，使用Dormand--Prince自适应积分\cite{dormand1980,shampine1997}作两档复核：最大步长与相对容差分别为$(0.2\,\mathrm s,10^{-10})$和$(0.1\,\mathrm s,10^{-11})$。分量绝对容差随数据保存。细化差值检查数值敏感性，不视为严格误差上界。

预算在积分前检查；自适应积分检测燃尽与落地，燃尽后切换至无推力运动，落地则终止。验收要求高度误差不超过500\,m、径向及切向速度误差各不超过10\,m/s，同时满足燃料预算且未触地。验收容差不等于搜索目标归一化尺度。

正式搜索共用16线程编译积分内核，MATLAB负责最终复核及MPC。另实现RTX 5060双精度CUDA核，与多核实现对同一96个候选的最大终态分量差约为$1.9\times10^{-9}$、目标最大差约为$8.4\times10^{-12}$。本机此小批量试算中多核实现更快，故正式比较统一用多核后端。计时排除编译与预热，不外推为其他硬件或批量下的性能排序。
